\section{هدف و مرز ماژول}
این ماژول فقط مسئول نگه‌داری داده است. از شبکه و
\lr{MQTT}
اطلاعی ندارد و تمام دستورهای
\lr{SQL}
در فایل زیر باقی می‌مانند:

\begin{LTR}
\begin{lstlisting}
code/collector/database.py
\end{lstlisting}
\end{LTR}

فایل پایگاه داده هنگام اجرا در مسیر زیر ساخته می‌شود و وارد
\lr{Git}
نمی‌شود:

\begin{LTR}
\begin{lstlisting}
runtime/iot.db
\end{lstlisting}
\end{LTR}

\section{چرا دو جدول؟}
هر پیام ممکن است چند
\lr{reading}
داشته باشد. برای جلوگیری از تکرار مشخصات پیام، داده در دو جدول نگه‌داری می‌شود:

\begin{itemize}
  \item
  \lr{telemetry\_batches}
  یک ردیف برای خود بسته، زمان دریافت، دستگاه و متن اصلی
  \lr{JSON}
  دارد؛
  \item
  \lr{sensor\_readings}
  برای هر سنسور یک ردیف عددی و قابل گزارش دارد.
\end{itemize}

اگر یک بسته شامل دما و رطوبت باشد، یک ردیف در جدول بسته‌ها و دو ردیف در جدول
\lr{reading}
ساخته می‌شود. فیلد
\lr{batch\_id}
این ردیف‌ها را به بستهٔ اصلی متصل می‌کند.

\section{ساختار جدول بسته‌ها}
ستون‌های اصلی جدول
\lr{telemetry\_batches}
عبارت‌اند از:

\begin{itemize}
  \item
  \lr{id}
  شناسهٔ داخلی و افزایشی پایگاه داده؛
  \item
  \lr{received\_at\_utc}
  زمان دریافت پیام در رایانه؛
  \item
  \lr{topic}
  موضوع واقعی
  \lr{MQTT}
  ؛
  \item
  \lr{device\_id}
  و
  \lr{message\_id}
  برای شناسایی منبع و نوبت نمونه‌برداری؛
  \item
  \lr{batch\_index}
  و
  \lr{batch\_count}
  برای داده‌هایی که در چند بسته ارسال شده‌اند؛
  \item
  \lr{payload\_json}
  نسخهٔ اصلی پیام برای عیب‌یابی آینده.
\end{itemize}

\section{ساختار جدول
\lr{reading}
ها}
جدول
\lr{sensor\_readings}
ستون‌های زیر را دارد:

\begin{LTR}
\begin{lstlisting}
id | batch_id | sensor_id | value | unit
\end{lstlisting}
\end{LTR}

قرارگرفتن مقدار در ستون عددی
\lr{REAL}
باعث می‌شود محاسبهٔ میانگین، کمینه و بیشینه ساده باشد. واحد جدا نگه‌داری می‌شود تا نام سنسور ثابت بماند.

\section{راه‌اندازی مستقل پایگاه داده}
این کد به
\lr{Python 3.10}
یا نسخهٔ جدیدتر نیاز دارد. نسخه را بررسی کنید:

\begin{LTR}
\begin{lstlisting}
python --version
\end{lstlisting}
\end{LTR}

از ریشهٔ پروژه محیط مجازی بسازید. این کار فقط بار اول لازم است:

\begin{LTR}
\begin{lstlisting}
# Windows PowerShell
python -m venv .venv
.\.venv\Scripts\Activate.ps1

# Linux / macOS
python3 -m venv .venv
source .venv/bin/activate
\end{lstlisting}
\end{LTR}

سپس وارد پوشهٔ
\lr{collector}
شوید و وابستگی را نصب کنید:

\begin{LTR}
\begin{lstlisting}
cd code/collector
python -m pip install -r requirements.txt
\end{lstlisting}
\end{LTR}

جدول‌ها را بدون اجرای کارگزار یا برد بسازید:

\begin{LTR}
\begin{lstlisting}
python database.py
\end{lstlisting}
\end{LTR}

خروجی مورد انتظار:

\begin{LTR}
\begin{lstlisting}
Database ready: .../runtime/iot.db
\end{lstlisting}
\end{LTR}

اجرای دوبارهٔ همین فرمان داده‌های قبلی را حذف نمی‌کند؛ فقط جدول‌ها و
\lr{index}
های وجودنداشته را می‌سازد.

\section{استفاده از ماژول در کد دیگر}
نمونهٔ زیر نشان می‌دهد که ماژول بدون
\lr{MQTT}
هم قابل استفاده است:

\begin{LTR}
\begin{lstlisting}[language=Python]
from database import SensorDatabase

payload = {
    "device_id": "esp8266-01",
    "message_id": 1,
    "sampled_at_ms": 5000,
    "batch_index": 1,
    "batch_count": 1,
    "readings": [
        {"sensor_id": "temperature", "value": (*@\lr{24.6}@*), "unit": "celsius"}
    ],
}

with SensorDatabase() as database:
    batch_id = database.save_telemetry_batch(
        "iot/devices/esp8266-01/telemetry",
        payload,
    )
    print(batch_id)
\end{lstlisting}
\end{LTR}

محیط
\lr{with}
تضمین می‌کند فایل پس از استفاده بسته شود. تابع ذخیره نیز بسته و
\lr{reading}
های آن را در یک
\lr{transaction}
می‌نویسد؛ بنابراین خطا در میانهٔ کار، بستهٔ ناقص باقی نمی‌گذارد.

\section{مشاهدهٔ داده‌های ذخیره‌شده}
برای مشاهدهٔ ده
\lr{reading}
آخر اجرا کنید:

\begin{LTR}
\begin{lstlisting}
python database.py --recent 10
\end{lstlisting}
\end{LTR}

نمونهٔ خروجی:

\begin{LTR}
\begin{lstlisting}
2026-08-22 10:13:51 esp8266-01 humidity=(*@\lr{58.2}@*) percent
2026-08-22 10:13:51 esp8266-01 temperature=(*@\lr{24.6}@*) celsius
\end{lstlisting}
\end{LTR}

برای بررسی مستقیم با ابزار خط فرمان
\lr{SQLite}
نیز می‌توان اجرا کرد:

\begin{LTR}
\begin{lstlisting}[language=SQL]
sqlite3 runtime/iot.db

.tables
.headers on
.mode column

SELECT sensor_id, value, unit
FROM sensor_readings
ORDER BY id DESC
LIMIT 10;
\end{lstlisting}
\end{LTR}

\section{نمونهٔ گزارش ساده}
میانگین هر سنسور با دستور زیر محاسبه می‌شود:

\begin{LTR}
\begin{lstlisting}[language=SQL]
SELECT
    sensor_id,
    COUNT(*) AS sample_count,
    ROUND(AVG(value), 2) AS average_value,
    MIN(value) AS minimum_value,
    MAX(value) AS maximum_value,
    unit
FROM sensor_readings
GROUP BY sensor_id, unit
ORDER BY sensor_id;
\end{lstlisting}
\end{LTR}

این
\lr{query}
شروع مناسبی برای مرحلهٔ گزارش است، اما هنوز رابط کاربری یا نمودار ساخته نشده است.

\section{پشتیبان‌گیری و پاک‌کردن دادهٔ آزمایشی}
برای پشتیبان‌گیری، ابتدا
\lr{collector}
را متوقف و سپس فایل
\lr{runtime/iot.db}
را کپی کنید. برای شروع دوباره از پایگاه خالی،
\lr{collector}
را متوقف کنید و فایل پایگاه داده را دستی حذف کنید؛ اجرای بعدی جدول‌ها را دوباره می‌سازد. این کار همهٔ داده‌های محلی را حذف می‌کند و نباید بدون نسخهٔ پشتیبان انجام شود.
